\chapter{Hạn chế và hướng phát triển}
\section{Hạn chế}
Hạn chế đầu tiên là nhãn sentiment được ánh xạ từ rating. Cách này nhanh và thực tế nhưng có thể nhiễu. Rating 3 sao không luôn là neutral; rating 5 sao vẫn có thể chứa phàn nàn; rating 1 sao có thể do yếu tố ngoài sản phẩm.

Hạn chế thứ hai là dataset sentiment chưa quá lớn. Với 5.162 review sau dedup, baseline có thể học tốt các từ khóa phổ biến nhưng chưa chắc tổng quát cho mọi ngành hàng. Lớp neutral ít hơn nên metric thấp hơn.

Hạn chế thứ ba là TF-IDF + SVM chưa hiểu ngữ cảnh sâu như PhoBERT/BERT. Mô hình có thể yếu với mỉa mai, phủ định phức tạp, câu dài nhiều ý và lỗi chính tả.

Hạn chế thứ tư là retrieval precision còn ở mức baseline. Manual Precision@5 khoảng 0,5067 cho thấy hệ thống đã truy xuất được evidence liên quan ở mức nhất định nhưng còn nhiều chỗ cải thiện. Top-k có thể lấy review cùng chủ đề nhưng sai intent hoặc sai polarity.

Hạn chế thứ năm là RAG chưa dùng LLM thật. Template-based generation an toàn và ổn định, nhưng khả năng diễn đạt và tổng hợp tự nhiên còn hạn chế. Nếu evidence phức tạp, template có thể chưa diễn giải tốt.

Hạn chế thứ sáu là Module 6 mới là MVP. Rule-based aspect sentiment chưa xử lý tốt phủ định, sarcasm, typo. Benchmark nhỏ chưa đủ để kết luận tổng quát.

\section{Hướng phát triển}
Hướng phát triển đầu tiên là fine-tune PhoBERT hoặc XLM-R cho sentiment classification. Điều này có thể cải thiện khả năng hiểu ngữ cảnh, đặc biệt với neutral/mixed sentiment.

Hướng thứ hai là xây dựng bộ dữ liệu ABSA tiếng Việt cho review TMĐT. Thay vì chỉ gán nhãn sentiment tổng thể, mỗi review có thể được gán aspect như chất lượng, giao hàng, đóng gói, dịch vụ shop và sentiment theo từng aspect.

Hướng thứ ba là huấn luyện supervised reranker cho retrieval. Reranker có thể học từ dữ liệu relevance thủ công để phân biệt review thật sự trả lời query và review chỉ cùng chủ đề.

Hướng thứ tư là mở rộng manual relevance evaluation. Bộ 15 query hiện tại quá nhỏ để kết luận mạnh. Cần nhiều query hơn, nhiều người chấm hơn và đo inter-annotator agreement.

Hướng thứ năm là tích hợp LLM có kiểm soát. Nếu dùng LLM, prompt phải bắt buộc dựa trên evidence, không được bịa, và câu trả lời phải có citation. LLM nên là lớp diễn đạt, không thay thế retrieval.

Hướng thứ sáu là cải thiện lọc spam/duplicate/noise. Có thể thêm model phát hiện review nhận xu, review không liên quan hoặc review copy-paste. Hướng thứ bảy là phát triển dashboard phân tích theo sản phẩm, category, thời gian và nhà bán.

\section{Định hướng production}
Nếu triển khai production, cần caching embedding/model, logging query, monitoring latency, kiểm soát version artifact, API authentication và cơ chế cập nhật index định kỳ. Hệ thống hiện tại mới là bản thử nghiệm học thuật end-to-end, chưa phải hệ thống production hoàn chỉnh.
